\documentclass[11pt]{article}

% ----------------------------------------------------------------------
% Step-by-step tutorial for the PGINN code: simulate,
% train and test the networks of one regime (10 loci, n = 10) with the
% paper's settings. The commands are in run_tutorial.sh; all numbers in
% this document are that script's output.
% ----------------------------------------------------------------------

\usepackage[margin=1in]{geometry}
\usepackage[T1]{fontenc}
\usepackage{lmodern}
\usepackage{array}
\usepackage{amsmath,amssymb}
\usepackage{booktabs}
\usepackage{graphicx}
\usepackage{listings}
\usepackage{xcolor}
\usepackage[hidelinks]{hyperref}

\lstset{basicstyle=\small\ttfamily, columns=fullflexible, keepspaces=true,
  frame=single, rulecolor=\color{black!30}, backgroundcolor=\color{black!3},
  breaklines=true, xleftmargin=0.5em, xrightmargin=0.5em, aboveskip=0.8em}
\lstdefinestyle{py}{language=Python, commentstyle=\color{black!55},
  keywordstyle=\bfseries}
\newcommand{\file}[1]{\texttt{#1}}
\usepackage{xspace}
\newcommand{\pginn}{PGINN\xspace}
\newcommand{\migrate}{\textsc{migrate}\xspace}

\title{\textbf{PGINN step by step:\\simulate, train and test the networks for one regime}}
\author{Peter Beerli}
\date{\today}

\emergencystretch=3em

\begin{document}
\maketitle

\begin{abstract}
The script \file{../code/reproduce.sh} rebuilds every figure and table of
the paper, but it calls more than 50 Python scripts, and the core of the
method is hard to see among them. This tutorial follows the core path
alone, for one regime: 10 loci per data set and $n{=}10$ samples. It
simulates genealogies under the sweep model, turns them into network
inputs, trains one network with the MSE loss and one with the \pginn
loss, and tests both on new simulated data. The commands are collected
in \file{run\_tutorial.sh}; the whole run takes about
70\,s on a laptop with 16 cores. The script uses the paper's priors,
sizes and random seeds, and its data and networks are identical, bit for
bit, to the paper's files for this regime (Section~\ref{sec:check}).
\end{abstract}

\tableofcontents

\section{Overview}\label{sec:overview}

\begin{table}[h]
\centering
\small
\begin{tabular}{@{}llll@{}}
\toprule
Step & Command & Script(s) & Output in \file{work/} \\
\midrule
1 & \file{genealogy} & \file{sweep.py}, \file{simtree.py} & \file{tp.txt}, \file{one.json}, \file{genealogy.pdf}\\
2 & \file{simulate}  & \file{simulate\_sweep.py} $\to$ \file{simtree.py} & \file{train-sweep/} and three more folders \\
3 & \file{features}  & \file{prepare.py} & \file{*-features.npy}, \file{*-params.npy}, \file{*-ages.npy} \\
4 & \file{train}     & \file{train\_pginn.py} & \file{model-\{mse,pginn\}.pt}, \file{scaler-*.npz} \\
5 & \file{test}      & \file{compare.py} & \file{recovery.pdf}, \file{recovery.csv} \\
6 & \file{score}     & \file{train\_pginn.predict()} & printed summary \\
-- & \file{check}    & -- & comparison with \file{../code/data}, \file{../code/model} \\
\bottomrule
\end{tabular}
\caption{The steps of \file{run\_tutorial.sh}. All scripts are in
\file{../code/}; paths in this document are relative to the folder \file{tutorial/}.}
\label{tab:steps}
\end{table}

Run all steps, or one step at a time, from any directory:
\begin{lstlisting}
cd tutorial                  # this folder; the scripts are in ../code
./run_tutorial.sh            # all steps, about 70 s
./run_tutorial.sh simulate   # one step; later steps need the earlier ones
./run_tutorial.sh check      # compare with the paper's data and networks
\end{lstlisting}
The script needs the Python packages of \file{../requirements.txt}
(\texttt{numpy}, \texttt{torch}, \texttt{matplotlib}); \texttt{msprime}
and \migrate are not needed. Everything is written to
\file{work/} (about 70\,MB).

\paragraph{The model in brief.}
A single population has the mutation-scaled size $\theta{=}c\mu$ and
the population-scaled selection coefficient $s{=}c\,s_g$.
A completed hard sweep shrinks the ancestral class of the sample toward
the past.
Toward the past, the frequency of the selected allele falls as
\begin{equation}
x(t;s) = \frac{x_0}{x_0 + (1-x_0)\,e^{st}}, \qquad x_0 = 0.999,
\label{eq:xtraj}
\end{equation}
down to the floor $x_{\min}{=}10^{-4}$.
Genealogies measure time in expected mutations per site, $u{=}\theta t$.
On this scale the size of the selected class and the coalescent rate for
$k$ lineages are
\begin{equation}
\theta_{\text{e}}(u) = \theta\,x(u/\theta;s), \qquad
\lambda_k(u) = \frac{k(k-1)}{\theta_{\text{e}}(u)} .
\label{eq:rate}
\end{equation}
The network estimates $(\theta,s)$ from the genealogies of a data set.

\section{Step 1: one genealogy by hand}\label{sec:step1}

The simulator \file{simtree.py} draws coalescent genealogies with a
piecewise-constant $\theta$.
A sweep enters only through a \emph{timepoints} file: each line holds an
age $u$ and the value of $\theta_{\text{e}}$ from that age on.
\file{sweep.py} writes this file. It places 40 breakpoints evenly in
$\log x$ between $0.999\,x_0$ and $x_{\min}$ and converts each to its age
by inverting Eq.~\ref{eq:xtraj}:
\begin{lstlisting}[style=py, title={\file{../code/sweep.py}, \texttt{make\_timepoints\_file} (core lines)}]
x_points = np.geomspace(x0 * 0.999, x_min, n_segments)
t_points = t_of_x(x_points, rate_per_age(s, theta), x0)
with open(path, "w") as f:
    for t, x in zip(t_points, x_points):
        f.write(f"{t:.10g} {theta * x:.10g}\n")
\end{lstlisting}
\texttt{rate\_per\_age(s, theta)} is $s/\theta$, so the trajectory is
evaluated directly on the age scale $u$.
Step~1 of the script writes the schedule for $\theta{=}0.01$, $s{=}100$
and simulates one genealogy with it and one neutral genealogy:
\begin{lstlisting}
python -c "import sweep; sweep.make_timepoints_file('tp.txt', 0.01, 100.0)"
python simtree.py -t 0.01 -i 10 -a 1.0 -sel 100 -tp tp.txt -l 1 --json -f one.json --seed 1
python simtree.py -t 0.01 -i 10 -a 1.0 -l 1 --json -f one-neutral.json --seed 1
\end{lstlisting}
The options are: \texttt{-t} $\theta$, \texttt{-i} sample size,
\texttt{-a 1.0} the Kingman coalescent (the simulator also has a
fractional coalescent, $\alpha{<}1$; the default \texttt{"1.0,1.0"}
would set up two populations), \texttt{-l} number of loci,
\texttt{--json} the output format used below, and \texttt{--seed} the
random seed.
The value of \texttt{-sel} is only recorded in the output; the sweep
itself acts through \texttt{-tp}.
The first lines of \file{tp.txt} are
\begin{lstlisting}
6.936475559e-05 0.00998001
0.0005593159182 0.007881142072
0.0006407142897 0.006223681175
\end{lstlisting}
Figure~\ref{fig:genealogy}a shows the schedule against the continuous
$\theta_{\text{e}}(u)$; Figure~\ref{fig:genealogy}b shows the two
genealogies.
The sweep genealogy has its root at $u{=}0.0012$, where
$\theta_{\text{e}}$ has collapsed; the neutral genealogy has its root at
$u{=}0.0066$.

\begin{figure}[t]
\centering
\includegraphics[width=\textwidth]{work/genealogy.pdf}
\caption{(a) $\theta_{\text{e}}(u)$ at $\theta{=}0.01$ for three values
of $s$ (Eq.~\ref{eq:rate}), and the 40-step schedule that
\file{simtree.py} receives for $s{=}100$ (dotted). (b) Number of
lineages against age for the two genealogies of step~1 (same random
seed). Made by \file{genealogy\_fig.py}.}
\label{fig:genealogy}
\end{figure}

\paragraph{The JSON output.}
\file{one.json} is a list with one entry per locus. Each entry has the
fields
\begin{center}
\small
\begin{tabular}{@{}llp{0.66\textwidth}@{}}
\toprule
Field & Length ($n{=}10$) & Content \\
\midrule
\texttt{params} & -- & $[[\theta],[\alpha],[s]]$ \\
\texttt{tree}   & -- & the genealogy in Newick format; tips \texttt{0\_0} \ldots\ \texttt{0\_9} (population\_tip) \\
\texttt{ages}   & 9  & ages of the $n{-}1$ internal nodes, oldest first \\
\texttt{age}    & 45 & age of the most recent common ancestor (MRCA) of each pair of tips \\
\texttt{inter}  & 45 & populations of the two tips of each pair (always \texttt{[0,0]} here) \\
\texttt{M}      & 55 & per pair: number of edges from the root to the MRCA; per tip: 1 \\
\texttt{m}      & 55 & per pair: length of the path from the root to the MRCA; per tip: its branch length \\
\bottomrule
\end{tabular}
\end{center}
The vectors \texttt{M} and \texttt{m} are the Kendall--Colijn tree
vectors, computed in \file{../code/Mm.py}.
Note that the JSON names are swapped relative to \file{Mm.py}, where
\texttt{m} is the edge count and \texttt{M} the path length; the JSON
field \texttt{m} holds the lengths.
Pairs are ordered by tip name, tips are listed after the pairs.

\section{Step 2: training, validation and test data}\label{sec:step2}

\file{simulate\_sweep.py} draws $(\theta,s)$ from the prior and calls
\file{simtree.py} once per draw, in parallel on all cores.
For $s{>}0$ it writes a timepoints file as in step~1.
Each draw has its own random seed, the base seed plus the draw number,
so a data set is reproducible exactly.
\begin{lstlisting}
PRIOR=(--theta-min 0.001 --theta-max 0.05 --theta-prior loguniform --s-prior loguniform --s-min 1)
python simulate_sweep.py -f train-sweep     -n 1000 -l 10 --nind 10 --smax 1000 --seed 100000 $PRIOR --s-zero-frac 0.25
python simulate_sweep.py -f test-sweep      -n 100  -l 10 --nind 10 --smax 1000 --seed 110000 $PRIOR --s-zero-frac 0.25
python simulate_sweep.py -f predict-sweep   -n 400  -l 10 --nind 10 --smax 1000 --seed 120000 $PRIOR
python simulate_sweep.py -f predict-neutral -n 400  -l 10 --nind 10 --smax 0    --seed 130000 $PRIOR
\end{lstlisting}
The four sets are
\begin{center}
\small
\begin{tabular}{@{}llp{0.33\textwidth}l@{}}
\toprule
Folder & Draws & Prior on $s$ & Use \\
\midrule
\file{train-sweep}     & 1000 & 0 with prob.\ 0.25, else log-uniform on $[1,1000]$ & training \\
\file{test-sweep}      & 100  & as training & validation (choice of the epoch) \\
\file{predict-sweep}   & 400  & log-uniform on $[1,1000]$ & test of recovery \\
\file{predict-neutral} & 400  & $s{=}0$ & test under neutrality \\
\bottomrule
\end{tabular}
\end{center}
In all four sets $\theta$ is log-uniform on $[0.001,0.05]$.
The file names follow \file{../code/reproduce.sh}, where the validation set
is called \file{test-sweep}.
Each draw is one data set of 10 independent loci (\texttt{-l 10}),
stored as one JSON file, for example
\file{train-sweep/test-000001-theta-0.002846-s-3.7509.json}.
Each folder also has a \file{manifest.csv} with the parameters and seed
of every draw.
In the training set 262 of the 1000 draws have $s{=}0$.

\section{Step 3: from genealogies to network inputs}\label{sec:step3}

\file{prepare.py} reads every JSON file of a folder and writes three
arrays, one row per draw (the counts are for $n{=}10$):
\begin{lstlisting}
python prepare.py --train train-sweep --interest 0,2 --prefix train-sweep
#  train-sweep-features.npy (1000, 225)   network input
#  train-sweep-params.npy   (1000, 2)     targets: [log theta, s]
#  train-sweep-ages.npy     (1000, 9)     mean internal-node ages, for the PGINN loss
\end{lstlisting}
\texttt{--interest 0,2} selects entries 0 and 2 of the flattened
\texttt{params}, that is $\theta$ and $s$.
The program prints ``Saved 3 averaged samples''; the 3 is the number of
arrays, not the number of draws.

\paragraph{The feature vector.}
Within each genealogy, the pair entries and the tip entries are sorted
by age, so that the vectors of different genealogies can be averaged
(tip labels have no meaning across independent loci).
The feature vector is then the mean over the 10 loci of four blocks.
With $P{=}n(n{-}1)/2+n{=}55$ entries per pair-and-tip block:
\begin{center}
\small
\begin{tabular}{@{}llll@{}}
\toprule
Columns & Block & Content & Scale in the network \\
\midrule
0--54    & 0 & path lengths, root to MRCA (JSON \texttt{m}) & log \\
55--109  & 1 & edge counts, root to MRCA (JSON \texttt{M}) & linear \\
110--164 & 2 & pair MRCA ages, then the tip branch lengths & log \\
165--224 & 3 & internal-node ages, youngest first, padded to $\max(20,n)$ rows $\times$ 3 & log (age column) \\
\bottomrule
\end{tabular}
\end{center}
The rows of block~3 are $(\text{age},\text{pop}_1,\text{pop}_2)$; for one
population the last two columns are zero, and rows beyond the $n{-}1$
internal nodes are zero padding.
In the paper's main networks the age-valued entries are replaced by
$\log(x+10^{-5})$. This transformation is the function \texttt{featlog}
in \file{../code/feature\_ages.py}.
\file{train\_pginn.py} applies it when called with
\texttt{--input-transform featlog} and stores the choice in the scaler
file, so the raw \file{-features.npy} serve as input throughout.
(\file{reproduce.sh} also runs \file{feature\_ages.py}, which writes
\file{-featlog.npy} and other input variants to disk for the comparisons
in the supplement; the tutorial does not need them.)

\paragraph{Targets and ages.}
\file{-params.npy} holds $[\log\theta, s]$ with $s$ on its natural
scale; \file{train\_pginn.py} turns the second column into
$\log(1{+}s)$.
\file{-ages.npy} holds, for every draw, the internal-node ages of each
locus sorted from the oldest to the youngest, averaged over the loci.
The \pginn loss evaluates the coalescent rates at these ages.

\section{Step 4: training}\label{sec:step4}

\begin{lstlisting}
b=(--x-train train-sweep-features.npy --y-train train-sweep-params.npy --ages-train train-sweep-ages.npy
   --x-val test-sweep-features.npy --y-val test-sweep-params.npy --ages-val test-sweep-ages.npy
   --input-transform featlog --seed 67)
python train_pginn.py $b --pginnratio 0   --out-model model-mse.pt   --out-scaler scaler-mse.npz
python train_pginn.py $b --pginnratio 0.1 --out-model model-pginn.pt --out-scaler scaler-pginn.npz
\end{lstlisting}

\paragraph{Network.}
A multilayer perceptron with two hidden layers of 256 units, each linear
layer followed by ReLU and batch normalization, and a linear output
layer with two values (\texttt{MLPRegressor}).
Inputs and targets are standardized with the means and standard
deviations of the training set.
Training uses AdamW (learning rate $10^{-3}$, weight decay $10^{-4}$),
batches of 64 and 300 epochs.

\paragraph{Loss.}
The network predicts $\hat\phi{=}(\log\hat\theta,\log(1{+}\hat s))$.
The training loss mixes the mean squared error with the \pginn penalty,
with weight $w$ (\texttt{--pginnratio}):
\begin{equation}
\mathcal{L} = (1-w)\,\text{MSE} + w\,\mathcal{L}_{\text{local}},
\qquad
\mathcal{L}_{\text{local}} =
\frac{1}{K}\sum_{k=2}^{n}
\left[
\log\frac{\hat\lambda_k(u_{(k)})}{\lambda_k(u_{(k)})} + \frac{\lambda_k(u_{(k)})}{\hat\lambda_k(u_{(k)})} - 1
\right],
\label{eq:loss}
\end{equation}
with $K{=}n{-}1$ and $u_{(k)}$ the age at which $k$ lineages merge into
$k{-}1$.
$w{=}0$ is the MSE network, and $w{=}0.1$ the \pginn network of the paper.
The penalty is the KL divergence between two exponential distributions
with the rates under the predicted and the true parameters, each
evaluated at the observed event age. In the code:
\begin{lstlisting}[style=py, title={\file{../code/train\_pginn.py}, \texttt{coalescent\_kl\_selection} (core lines)}]
K = ages.shape[1]
k = torch.arange(2, K + 2, device=ages.device, dtype=ages.dtype)  # [2,3,...,K+1]
k_km1 = k * (k - 1)  # (K,)

x_hat = x_trajectory_torch(ages, (s_hat / theta_hat).unsqueeze(1))     # (B,K)
x_true = x_trajectory_torch(ages, (s_true / theta_true).unsqueeze(1))  # (B,K)

theta_eff_hat = theta_hat.unsqueeze(1) * x_hat + eps
theta_eff_true = theta_true.unsqueeze(1) * x_true + eps

lam_hat = (k_km1.unsqueeze(0) / theta_eff_hat).clamp(max=MAX_LAMBDA)
lam = (k_km1.unsqueeze(0) / theta_eff_true).clamp(max=MAX_LAMBDA)

kl_per_interval = torch.log(lam_hat / lam + eps) + lam / lam_hat - 1.0
return kl_per_interval.mean(dim=1)
\end{lstlisting}
Before these lines the function undoes the standardization of the
predicted and the true targets and converts them to $\hat\theta,\hat s$
and $\theta,s$.
Column 0 of \texttt{ages} is the root, so it belongs to $k{=}2$.
The mixing of the two terms is in \texttt{train\_epoch}:
\begin{lstlisting}[style=py, title={\file{../code/train\_pginn.py}, \texttt{train\_epoch}}]
pred = model(xb)
mse = mse_criterion(pred, yb) if s_weight == 1.0 else weighted_mse(pred, yb, s_weight)
kl_loss = term(pred, yb, y_mean, y_std, ages).mean()
loss = (1.0 - pginnratio) * mse + pginnratio * kl_loss
\end{lstlisting}
The option \texttt{--loss} selects the other penalties of the paper
(\texttt{exact}, \texttt{fisher}, \texttt{process},
\texttt{process-fwd}); they need other age files
(Section~\ref{sec:further}).

\paragraph{Choice of the epoch and the output.}
After every epoch the loss on the validation set is computed; the
network of the epoch with the lowest validation loss is kept.
The last lines of the two training logs (\file{work/train-*.log}) are
\begin{lstlisting}
==> work/train-mse.log <==
Best validation loss: 0.06665127277374268 at epoch 5
==> work/train-pginn.log <==
Best validation loss: 0.07293154954910279 at epoch 300
\end{lstlisting}
For this training run the MSE network is the one of epoch~5: its
validation loss fluctuates strongly in the first epochs (0.075, 0.081,
0.067, 0.097, 0.108 in epochs 3--7) and does not fall below the value of
epoch~5 later; at epoch 300 it is 0.080.
Each run writes two files: the network weights (\file{.pt}) and a scaler
file (\file{.npz}) with the standardization constants, the
$\log(1{+}s)$ link and the input transformation.
Both files are needed for prediction.
The weights are reproducible bit for bit only with the same random seed
and the same number of CPU threads (\texttt{--threads}, default 2).

\section{Step 5: test of recovery}\label{sec:step5}

\begin{lstlisting}
python compare.py --x predict-sweep-features.npy --y predict-sweep-params.npy \
  --model-mse model-mse.pt --scaler-mse scaler-mse.npz \
  --model-pginn model-pginn.pt --scaler-pginn scaler-pginn.npz \
  --out-fig recovery.pdf --out-table recovery.csv
\end{lstlisting}
\file{compare.py} estimates $(\theta,s)$ for the 400 test data sets
with both networks and reports bias, root mean squared error and $R^2$
(Table~\ref{tab:recovery}, Figure~\ref{fig:recovery}).
$s$ is evaluated on the $\log(1{+}s)$ scale, the scale of its prior and
of the network's target.
\begin{table}[h]
\centering
\small
\begin{tabular}{@{}llrrr@{}}
\toprule
Parameter & Loss & Bias & RMSE & $R^2$ \\
\midrule
$\theta$       & MSE   & 0.00041  & 0.0030 & 0.939 \\
               & \pginn & $-0.00046$ & 0.0028 & 0.946 \\
$\log(1{+}s)$  & MSE   & $-0.338$ & 0.622 & 0.881 \\
               & \pginn & $-0.290$ & 0.671 & 0.862 \\
\bottomrule
\end{tabular}
\caption{Recovery on the 400 test data sets (\file{work/recovery.csv},
rounded). The file also has the rows for $s$ on its natural scale.}
\label{tab:recovery}
\end{table}
\file{reproduce.sh} writes the same table to
\file{../figures/recovery\_10l\_featlog.csv}.
It is one training run on one data set. The comparison of the losses in
\file{nested\_eval.py} uses 10 independent data sets with three
training runs each.

\begin{figure}[t]
\centering
\includegraphics[width=\textwidth]{work/recovery.pdf}
\caption{Estimated against true $\theta$ and $\log(1{+}s)$ for the 400
test data sets, MSE and \pginn network (\file{compare.py}).}
\label{fig:recovery}
\end{figure}

\section{Step 6: estimates for new data}\label{sec:step6}

New data enter the networks as feature vectors with the layout of
step~3. The function \texttt{predict} in \file{train\_pginn.py} applies
the scaler, runs the network and undoes the links:
\begin{lstlisting}[style=py, title={\file{../code/train\_pginn.py}, \texttt{predict} (core lines)}]
scaler = np.load(scaler_path, allow_pickle=True)
if "input_transform" in scaler and str(scaler["input_transform"]) == "featlog":
    from feature_ages import featlog
    X = featlog(X)
x_std = (X - scaler["x_mean"]) / scaler["x_std"]
model = MLPRegressor(X.shape[1], scaler["y_mean"].shape[0])
model.load_state_dict(torch.load(model_path, map_location="cpu"))
model.eval()
with torch.no_grad():
    z = model(torch.from_numpy(x_std).float()).numpy() * scaler["y_std"] + scaler["y_mean"]
theta = np.clip(np.exp(z[:, 0]) + MIN_THETA, *THETA_SUPPORT)
s = np.clip(np.expm1(z[:, 1]), *S_SUPPORT)
\end{lstlisting}
Estimates are bounded to the support of the training prior,
$\theta\in[0.001,0.05]$ and $s\in[0,1000]$.
Step~6 scores the 400 neutral test data sets:
\begin{lstlisting}
mse    median s_hat   1.31   s_hat < 5: 0.96   corr(log theta_hat, log theta) 0.992
pginn  median s_hat   1.06   s_hat < 5: 0.95   corr(log theta_hat, log theta) 0.993
\end{lstlisting}
Under neutrality both networks give small $\hat s$ for most data sets.
The paper uses $\hat s$ as the detection statistic, with a threshold
from the neutral sets of another data set (\file{nested\_eval.py}).
For real sequence data the genealogies are not known; the paper scores
the posterior genealogies of \migrate
(\file{real\_adh\_migrate\_scan.py}); that script writes each tree in
the JSON format of step~1 (with \file{Mm.py}) and builds the feature
vectors with the code of \file{prepare.py}.

\section{Check against the paper's files}\label{sec:check}

\file{./run\_tutorial.sh check} compares the arrays of step~3 and the
networks of step~4 with the paper's files:
\begin{lstlisting}
work/{set}-{features,params,ages}.npy  vs  ../code/data/{set}-10l-{features,params,ages}.npy
work/model-mse.pt                      vs  ../code/model/model-mse-10l-featlog.pt
work/model-pginn.pt                    vs  ../code/model/model-pginn-10l-0.1-featlog.pt
\end{lstlisting}
The networks are those of training run 67, created by the \texttt{simulate} and
\texttt{revision-train} stages of \file{reproduce.sh} (or fetched from
Zenodo with \file{fetch\_data.sh}).
On 2026-10-23 all twelve arrays and both networks were identical.
If they differ on another machine, the most likely reasons are a
different number of CPU threads or a different version of
\texttt{torch} or \texttt{numpy}.

\section{Further: the rest of the code}\label{sec:further}

\paragraph{Other regimes.}
\file{reproduce.sh} runs steps 2--4 for six regimes, with the same
commands and other values of \texttt{-l}, \texttt{--nind} and the seeds:
\begin{center}
\small
\begin{tabular}{@{}lrrr@{}}
\toprule
Regime & Loci & $n$ & Base seed \\
\midrule
\texttt{10l}     & 10 & 10  & 100000 \\
\texttt{1l}      & 1  & 10  & 200000 \\
\texttt{10l-n20} & 10 & 20  & 300000 \\
\texttt{1l-n20}  & 1  & 20  & 400000 \\
\texttt{1l-n176} & 1  & 176 & 600000 \\
\texttt{1l-n100} & 1  & 100 & 800000 \\
\bottomrule
\end{tabular}
\end{center}
The validation, selected test and neutral test sets use the base seed
plus 10000, 20000 and 30000.
The paper trains 10 runs per regime and loss (seeds 67--76).
The regime \texttt{1l-n176} is used for the \emph{Adh} data and
\texttt{1l-n100} for the simulated chromosomes (stage
\texttt{chrom-simulate}).

\paragraph{Other penalties.}
The penalties of the paper's penalty family differ in the option
\texttt{--loss} and in the file passed as \texttt{--ages-*}:
\begin{center}
\small
\begin{tabular}{@{}lll@{}}
\toprule
Penalty (paper) & \texttt{--loss} & \texttt{--ages-train}, \texttt{--ages-val} \\
\midrule
local rate divergence      & \texttt{local} (default) & \file{-ages.npy} \\
exact log-likelihood ratio & \texttt{exact}       & \file{-locusages.npy} (\file{per\_locus\_ages.py}) \\
process divergences        & \texttt{process}, \texttt{process-fwd} & \file{-locusages.npy} \\
Fisher penalty             & \texttt{fisher}      & \file{-fisher.npy} (\file{fisher\_info.py}) \\
weighted MSE               & \texttt{--pginnratio 0 --s-weight 5} & \file{-ages.npy} (not used) \\
\bottomrule
\end{tabular}
\end{center}

\paragraph{Map of the scripts.}
The other scripts of \file{../code/} fall into these groups; the file
\file{../reproducibility.txt} lists which script writes which output
file.
\begin{center}
\small
\begin{tabular}{@{}>{\raggedright\arraybackslash}p{0.25\textwidth}>{\raggedright\arraybackslash}p{0.71\textwidth}@{}}
\toprule
Group & Scripts \\
\midrule
Core (this tutorial) & \file{sweep.py}, \file{simtree.py} (with \file{tree.py}, \file{Mm.py}), \file{simulate\_sweep.py}, \file{prepare.py}, \file{feature\_ages.py}, \file{train\_pginn.py}, \file{compare.py} \\
Inputs for other penalties & \file{per\_locus\_ages.py}, \file{fisher\_info.py} \\
Choice of the weights & \file{lambda\_sweep\_val.py}, \file{variant\_sweep\_val.py}, \file{w1\_diagnostic.py}, \file{w1\_check.py} \\
Comparison of the losses & \file{nested\_eval.py} (main table), \file{variance\_mechanism.py}, \file{combined\_detector.py}, \file{eval\_n176.py}, \file{penalty\_share.py}, \file{loss\_timing.py}, \file{asymmetry\_fig.py} \\
Baselines & \file{mle\_baseline.py}, \file{mle\_check.py}, \file{ages\_regression.py}, \file{star\_ratio\_appendix.py} \\
Checks of the model and the penalties & \file{calibrate\_sweep.py}, \file{discretization\_check.py}, \file{fisher\_check.py}, \file{process\_kl\_check.py}, \file{neutral\_limit.py} \\
\texttt{msprime} chromosomes & \file{msprime\_scan*.py}, \file{msprime\_metrics.py}, \file{robustness\_demography.py} \\
Chromosomes through \migrate & \file{chrom\_pipeline.sh}, \file{chrom\_sweep\_sim.py}, \file{chrom\_*.py}, \file{average\_runs.py} \\
\emph{Adh} data & \file{four\_gamete\_test.py}, \file{make\_migrate\_windows.py}, \file{extract\_dgrp\_region.py}, \file{real\_adh\_migrate\_scan.py}, \file{adh\_*.py}, \file{real\_data\_summary\_fig.py}, \file{mcdonald\_kreitman\_test.py} \\
Other & \file{detect\_selection.py}, \file{paired\_test.py}, \file{../talk/growth\_scan\_talk.py} (growth scan) \\
\bottomrule
\end{tabular}
\end{center}

\end{document}
